% ===================================================================
%  ТАБЛИЦА № 9 — Горы и курорт. — Лес и охота.
%  Traduction 2026 d'apres texto/io, controlee sur texto/fr, texto/en
%  et texto/es. Consigne : voir texto/ru/10-tablica-01.tex.
% ===================================================================

%%K t09-apar-1 apar
\VUsaut{4.0mm}
\VUtitre{15.0pt}{60}{ТАБЛИЦА № 9}
\VUsaut{3.0mm}

%%K t09-tit-1 sub
\VUcentre{12.0pt}{0}{\textit{Первая сцена.}}
\VUsaut{3.0mm}

%%K t09-tit-2 sub
\VUpk{11.4pt}{40}{Горы. --- Курортный городок.}
\VUsaut{3.0mm}

%%K t09-c1-01-1 p
1. --- Я уже неделю в Пратце. Не могу передать всего изумления, какое я
испытал, став перед чудесной \VUgras{горной цепью} \textsuperscript{(1)} Пиренеев,
чей \VUgras{гребень} \textsuperscript{(2)} вырезается на синем небе, как
исполинская пила.

%%K t09-c1-02-1 p
2. --- Я не воображал, что пиренейские \VUgras{вершины} \textsuperscript{(3)}
могут достигать такой \VUgras{высоты}, и остался поражён великолепными
\VUgras{ледниками} \textsuperscript{(5)} и снежными \VUgras{пиками}
\textsuperscript{(4)}, по которым несутся такие страшные \VUgras{лавины}.

%%K t09-tit-3 sub
\VUsaut{3.0mm}
\VUpk{10.2pt}{40}{Горный вид.}
\VUsaut{3.0mm}

%%K t09-c1-03-1 p
3. --- На другой же день по приезде я пошёл к старому потухшему
\VUgras{вулкану} \textsuperscript{(6)} возле Пратца. На голом бесплодном краю
\VUgras{кратера} ещё ясно видны следы, оставленные \VUgras{лавой}, которая
текла несколько веков назад по \VUgras{склону горы} \textsuperscript{(7)}. На
одном из \VUgras{откосов} мне удалось найти несколько кусков этой лавы.

%%K t09-c1-04-1 p
4. --- Пратц стоит на границе, и \VUgras{крепость} \textsuperscript{(8)},
стерегущая \VUgras{перевал} \textsuperscript{(27)} между Францией и Испанией,
очень похожа на те старые замки по берегам Рейна, которые словно
высматривают \VUgras{добычу} с вершины своей скалы.

%%K t09-c1-05-1 p
5. --- Огромный \VUgras{орёл} \textsuperscript{(9)}, у которого гнездо недалеко
от \VUgras{форта} \textsuperscript{(8)}, часто привлекает моё внимание. Эта
\VUgras{хищная птица} с могучим \VUgras{клювом} часто приносит птенцам
ягнёнка или козлёнка, зажатого в \VUgras{когтях}. \VUgras{Крылья} у неё
так велики, что она долго может парить со своей тяжёлой ношей.

%%K t09-tit-4 sub
\VUsaut{3.0mm}
\VUpk{10.2pt}{40}{Пешеход.}
\VUsaut{3.0mm}

%%K t09-c1-06-1 p
6. --- Самая приятная здешняя прогулка --- поход к \VUgras{водопаду}
\textsuperscript{(10)} Ко. \VUgras{Падающая вода} рушится вихрем и потом исчезает
хорошеньким \VUgras{ручьём} в \VUgras{еловом лесу} \textsuperscript{(11)} к западу
от городка. Мы поднялись этим лесом до \VUgras{метеорологической станции}
\textsuperscript{(12)} и с вершины \VUgras{вышки} охватили взглядом всю
\VUgras{панораму} округи. \VUgras{Вид} великолепен, и \VUgras{природа},
кажется, расточила этому краю все сокровища, какие были в её
распоряжении.

%%K t09-c1-07-1 p
7. --- Чтобы спуститься, мы взяли \VUgras{фуникулёр} \textsuperscript{(13)} и
сошли у маленькой \VUgras{станции} возле \VUgras{развалин} \textsuperscript{(14)}
старого \VUgras{феодального замка}, в котором некогда жили владетели
Пратца.

%%K t09-c1-08-1 p
8. --- Бедный замок! Что должен он думать, видя внизу под собою опрятный
\VUgras{курортный городок} \textsuperscript{(15)}, ставший модным
\VUgras{водным курортом}?

%%K t09-c1-08-2 p
\VUgras{Климат} здесь так приятен, а \VUgras{минеральная вода} так
целебна, что \VUgras{курс лечения}, который проходят в \VUgras{санатории}
\textsuperscript{(17)}, --- истинное удовольствие.

%%K t09-tit-5 sub
\VUsaut{3.0mm}
\VUpk{10.2pt}{40}{Приятный курортный городок.}
\VUsaut{3.0mm}

%%K t09-c1-09-1 p
9. --- И правда, сделано всё, чтобы пребывание было приятным. Построили
большой \VUgras{виадук} \textsuperscript{(16)} для железной дороги;
\VUgras{парк} \textsuperscript{(18)} при \VUgras{курзале}, где дают
\VUgras{концерты}, разбит прелестно. Несколько изящных \VUgras{вилл} \textsuperscript{(19)}
выстроили вокруг \VUgras{казино} \textsuperscript{(21)}, а очень хорошо
содержимое \VUgras{шоссе} \textsuperscript{(22)} делает передвижение чрезвычайно
лёгким.

%%K t09-c1-10-1 p
10. --- Простая \VUgras{насыпь} \textsuperscript{(23)} отделяет \VUgras{дорогу}
\textsuperscript{(20)} от самой \VUgras{долины} \textsuperscript{(25)}, на дне которой
лежит грациозное \VUgras{озерцо} \textsuperscript{(26)}, питающее прозрачный
\VUgras{ручей} \textsuperscript{(24)}.

%%K t09-c1-11-1 p
11. --- По другую сторону долины разрабатывают железный \VUgras{рудник}
\textsuperscript{(28)}. Я несколько раз ходил смотреть, как \VUgras{рудокопы}
спускаются в \VUgras{шахту}, и даже входил в \VUgras{штольню}, где они
проводят день, добывая \VUgras{руду}, в которой заключён \VUgras{металл}.

%%K t09-c1-12-1 p
12. --- Возле рудника построили большой \VUgras{литейный завод}, чтобы
тут же употребить извлекаемые из него богатства. \VUgras{Домна}
\textsuperscript{(29)}, топимая здешним обильным \VUgras{углём}, позволяет
получать \VUgras{чугун} быстро и дёшево.

%%K t09-c1-13-1 p
13. --- Недалеко от рудника разрабатывают \VUgras{каменоломню}
\textsuperscript{(30)}. Старый \VUgras{каменолом} откалывает своей \VUgras{киркой}
не \VUgras{гранит}, не \VUgras{порфир} и даже не \VUgras{песчаник}, а простой
\VUgras{строительный камень} для постройки домов.

%%K t09-c1-14-1 p
14. --- Одно из лучших украшений этого края --- \VUgras{ель}
\textsuperscript{(31)}. Повсюду --- на дне \VUgras{оврага} \textsuperscript{(32)}, на
берегу \VUgras{потока} \textsuperscript{(33)}, у \VUgras{пропасти}
\textsuperscript{(34)} или обрыва --- её тонкие иглы вносят свою зелёную ноту.

%%K t09-tit-6 sub
\VUsaut{3.0mm}
\VUpk{10.2pt}{40}{Прогулки пешком.}
\VUsaut{3.0mm}

%%K t09-c1-15-1 p
15. --- Я пользуюсь каникулами, чтобы совершать долгие прогулки.
\VUgras{Тропы} \textsuperscript{(35)}, вьющиеся по горе, позволяют пройти
несколько \VUgras{километров}, а часто и несколько \VUgras{вёрст}, не
замечая расстояния.

%%K t09-c1-15-2 p
\VUgras{Верстовые столбы} \textsuperscript{(36)} и \VUgras{указатели}
\textsuperscript{(37)} размечают тропы, на которых часто встречаешь молодого
\VUgras{горца} \textsuperscript{(38)} с \VUgras{котомкой} \textsuperscript{(40)} на
боку, несущего \VUgras{сурка} \textsuperscript{(39)}, или какого-нибудь
\VUgras{цыгана} \textsuperscript{(41)}, у которого \VUgras{меховая шапка}
\textsuperscript{(42)} странно спорит со старым рваным \VUgras{плащом}
\textsuperscript{(43)}. Он ведёт на \VUgras{цепи} учёного \VUgras{медведя}
\textsuperscript{(44)}.

%%K t09-tit-7 sub
\VUpk{10.2pt}{40}{Восхождения.}
\VUsaut{3.0mm}

%%K t09-c1-16-1 p
16. --- Почти каждый день я хожу с \VUgras{проводником} \textsuperscript{(48)}
на \VUgras{восхождение}, и я очень горд, когда с \VUgras{ранцем}
\textsuperscript{(47)} за спиной и \VUgras{альпенштоком} в руке, как настоящий
\VUgras{турист} \textsuperscript{(46)}, беру приступом \VUgras{скалы}
\textsuperscript{(45)}.

%%K t09-c1-17-1 p
17. --- С вершины горы люди на равнине кажутся не больше муравьёв;
\VUgras{отара овец} \textsuperscript{(50)}, пасущаяся на \VUgras{пастбище}
\textsuperscript{(49)}, похожа на детскую игрушку. \VUgras{Сосна}
\textsuperscript{(51)} или даже деревянное \VUgras{шале} \textsuperscript{(52)} --- едва
заметная точка на зелени.

%%K t09-c1-18-1 p
18. --- На этих прогулках мы часто встречаем на \VUgras{тропе}
\textsuperscript{(53)} \VUgras{охотника на серн} \textsuperscript{(54)}, несущего на
плече только что убитого зверя.

%%K t09-c1-19-1 p
19. --- Я положил в свой гербарий очень примечательный лист
\VUgras{папоротника} \textsuperscript{(55)} и хорошенькую розовую веточку
\VUgras{вереска} \textsuperscript{(57)}, которую сорвал у входа в маленькую
\VUgras{пещеру} \textsuperscript{(56)} на последней прогулке.

%%K t09-tit-8 sub
\VUsaut{3.0mm}
\VUcentre{12.0pt}{0}{\textit{Вторая сцена.}}
\VUsaut{3.0mm}

%%K t09-tit-9 sub
\VUpk{11.4pt}{40}{Лес. --- Охота.}
\VUsaut{3.0mm}

%%K t09-c2-01-1 p
1. --- Вчера я провёл в лесу восхитительный день. Деятельная жизнь
зверей и \VUgras{насекомых} \textsuperscript{(5)}, которые там живут, очень меня
занимала.

%%K t09-c2-01-2 p
\VUgras{Паук} \textsuperscript{(1)}, ткущий свою \VUgras{паутину} \textsuperscript{(2)}
между двух веток, чтобы поймать пролетающую \VUgras{муху} \textsuperscript{(3)};
\VUgras{улитка} \textsuperscript{(4)}, тяжело тащащаяся со своим домиком; жук-олень,
высматривающий добычу; хлопотливый муравей, работающий возле
\VUgras{муравейника} \textsuperscript{(6)} --- все эти малые твари сновали и
трудились с необычайной силой.

%%K t09-c2-02-1 p
2. --- \VUgras{Змея} \textsuperscript{(7)}, которую я сперва принял за
\VUgras{гадюку}, а она оказалась безобидным \VUgras{ужом}, свернулась под
стволом дерева и время от времени высовывала свою тонкую головку, всегда
как будто готовую ужалить. Передо мною тяжело ползла большая
\VUgras{слизень} \textsuperscript{(8)}, съевшая одну из вкусных маленьких
\VUgras{земляник} \textsuperscript{(11)}, которых здесь множество.

%%K t09-c2-03-1 p
3. --- Земля была устлана \VUgras{лесными цветами}: \VUgras{первоцветы}
\textsuperscript{(9)}, \VUgras{барвинки} \textsuperscript{(10)} и много других
неизвестных мне растений цвели среди \VUgras{кустиков травы}
\textsuperscript{(12)}.

%%K t09-c2-04-1 p
4. --- В этом краю \VUgras{трюфель} почти так же обыкновенен, как
\VUgras{гриб} \textsuperscript{(14)}, и \VUgras{поросёнок} \textsuperscript{(13)},
принадлежащий одному леснику, жадно рылся, не найдётся ли их.

%%K t09-tit-10 sub
\VUsaut{3.0mm}
\VUpk{10.2pt}{40}{Браконьерство.}
\VUsaut{3.0mm}

%%K t09-c2-05-1 p
5. --- Я застал \VUgras{конец} сцены, которая меня очень позабавила.
Благодаря чутью своей \VUgras{таксы} \textsuperscript{(15)} \VUgras{лесник}
\textsuperscript{(16)} только что накрыл \VUgras{браконьера} \textsuperscript{(22)} в
ту минуту, когда тот собирался унести убитую им \VUgras{дичь}
\textsuperscript{(17)}. Предпочтя расстаться с добычей, лишь бы не быть
задержанным, браконьер убежал, и \VUgras{заяц} \textsuperscript{(18)},
\VUgras{лань} \textsuperscript{(19)}, \VUgras{косуля} \textsuperscript{(20)} и
\VUgras{кабан} \textsuperscript{(21)} остались лежать на траве, к великой радости
лесника.

%%K t09-c2-06-1 p
6. --- Поражает разнообразие деревьев в этом лесу. \VUgras{Буки}
\textsuperscript{(23)} и \VUgras{берёзы} \textsuperscript{(26)}, \VUgras{сосны}, дающие
\VUgras{смолу} \textsuperscript{(27)}, \VUgras{дубы} \textsuperscript{(29)} и много
других переплетают свои ветви.

%%K t09-c2-06-2 p
\VUgras{Плющ} \textsuperscript{(24)}, цепляющийся за стволы высоких деревьев,
придаёт \VUgras{высокоствольному лесу} \textsuperscript{(25)} яркую зелень,
которая приятно сливается с бледным и мягким оттенком \VUgras{мха}
\textsuperscript{(31)}, где \VUgras{зелёный дятел} \textsuperscript{(30)} ищет
насекомых.

%%K t09-tit-11 sub
\VUsaut{3.0mm}
\VUpk{10.2pt}{40}{Лесной вид.}
\VUsaut{3.0mm}

%%K t09-c2-07-1 p
7. --- Старые дубы меня очаровали. Их большие узловатые \VUgras{корни}
\textsuperscript{(32)}, наполовину вышедшие из земли, придают им великолепный вид
силы и мощи, и я не понимаю, как \VUgras{владелец} \textsuperscript{(35)} может
решиться их валить и отдавать \VUgras{пиле} \textsuperscript{(34)}
\VUgras{дровосека} \textsuperscript{(33)}. Бедные \VUgras{стволы}
\textsuperscript{(36)}, лишённые \VUgras{коры} \textsuperscript{(37)} или расколотые
\VUgras{топором} \textsuperscript{(38)} \VUgras{подёнщика} \textsuperscript{(39)},
печалят меня.

%%K t09-c2-08-1 p
8. --- Рядом старый \VUgras{угольщик} \textsuperscript{(41)} сложил своей
\VUgras{лопатой} \VUgras{кучу} \textsuperscript{(42)} угля, а старуха уносила
\VUgras{вязанку} \textsuperscript{(40)} \VUgras{хвороста} из веток срубленных
деревьев.

%%K t09-tit-12 sub
\VUsaut{3.0mm}
\VUpk{10.2pt}{40}{Охота.}
\VUsaut{3.0mm}

%%K t09-c2-09-1 p
9. --- Я собирался немного отдохнуть в \VUgras{доме лесника}
\textsuperscript{(46)}, но, дойдя до маленького \VUgras{пруда} \textsuperscript{(43)},
по которому плавал красивый \VUgras{лебедь} \textsuperscript{(45)}, я заметил, что
недавнее \VUgras{половодье} \textsuperscript{(44)} обратило тропу в настоящее
\VUgras{болото}. Пришлось сделать крюк, и, проходя мимо \VUgras{кедра}
\textsuperscript{(49)}, \VUgras{тополей} \textsuperscript{(48)} и \VUgras{вяза}
\textsuperscript{(47)}, что стоят на опушке, я увидел, как из \VUgras{чащи}
\textsuperscript{(54)} вырвался великолепный \VUgras{олень}
\textsuperscript{(50)}, преследуемый неутомимой \VUgras{сворой} \textsuperscript{(51)}, которую горячил
\VUgras{рог} \textsuperscript{(53)} \VUgras{конных охотников} \textsuperscript{(52)}.
Бедный зверь совсем выбился из сил. Он упал, издыхая, в нескольких шагах
от меня.

%%K t09-c2-10-1 p
10. --- Сын лесника, привлечённый шумом \VUgras{охоты}, вышел из дома, и мы
вдвоём вернулись в лес. Он приметил \VUgras{сороку} \textsuperscript{(56)} на
ветке старого дуба. Он думал, что гнездо её спрятано в \VUgras{листве}
\textsuperscript{(58)}, и хотел его взять.

%%K t09-c2-10-2 p
Проворный, как \VUgras{белка} \textsuperscript{(61)}, он лез с \VUgras{ветки}
\textsuperscript{(55)} на ветку, но ничего не нашёл и только вспугнул старую
\VUgras{ворону} \textsuperscript{(60)}, сидевшую на \VUgras{суку} \textsuperscript{(57)}.
\VUsaut{4.0mm}
\VUfilet{22mm}
